% LUCAS publication data -- auto-generated by tools/lucas_paper_data.py
\newcommand{\LucasPrimeTable}{%
\begin{table}[ht]
\centering
\caption{Lucas-prime $\phi$ periods computed by the publication oracle.}
\label{tab:lucas}
\footnotesize
\setlength{\tabcolsep}{5pt}
\begin{tabular}{rrr}
\toprule
$p$ & $L_p$ & Period \\
\midrule
5 & 11 & 10 \\
7 & 29 & 14 \\
11 & 199 & 22 \\
13 & 521 & 26 \\
17 & 3571 & 34 \\
19 & 9349 & 38 \\
\bottomrule
\end{tabular}
\end{table}
}

\newcommand{\LucasTimingTable}{%
\begin{table}[ht]
\centering
\caption{Accepted-operation latency of the reusable MAC core. PMUL's three
cycles are busy cycles after acceptance; PINV is data-dependent but bounded by
a 64-iteration watchdog. Transport and sidecar scheduling are excluded.}
\label{tab:timing}
\footnotesize
\setlength{\tabcolsep}{3pt}
\begin{tabular}{lll}
\toprule
Opcode & Core latency & Arithmetic path \\
\midrule
PSCALE & 1 cycle & shift/add, zero multiply \\
PCHIRAL & 1 cycle & add/subtract, zero multiply \\
PMUL & 3 busy cycles & four scalar products \\
PINV & bounded iterative & norm, BEEA, scalar products \\
PHSLK & 1 cycle & two cross-products \\
\bottomrule
\end{tabular}
\end{table}
}

\newcommand{\LucasOperationTable}{%
\begin{table}[ht]
\centering
\caption{Selected arithmetic vectors recomputed by the publication oracle.}
\label{tab:ops}
\footnotesize
\setlength{\tabcolsep}{3pt}
\begin{tabular}{lll}
\toprule
Operation & Oracle pair $(a,b)$ & Match \\
\midrule
$\phi(3+5\phi)$ & $(5,8)$ & $\checkmark$ \\
$\overline{3+5\phi}$ & $(8,516)$ & $\checkmark$ \\
$(3+5\phi)(2+7\phi)$ & $(41,66)$ & $\checkmark$ \\
$(3+5\phi)^{-1}$ & $(513,5)$ & $\checkmark$ \\
$(3+5\phi)(3+5\phi)^{-1}$ & $(1,0)$ & $\checkmark$ \\
\bottomrule
\end{tabular}
\end{table}
}

\newcommand{\LucasResourceTable}{%
\begin{table*}[t]
\centering
\caption{Mapped implementation artifacts. Backend-native logic counts are not
cross-vendor equivalents. ``Silicon'' denotes the bounded vectors recorded in
the hardware evidence ledger, not exhaustive input coverage.}
\label{tab:resources}
\footnotesize
\setlength{\tabcolsep}{5pt}
\begin{tabular}{llll}
\toprule
Artifact & Logic/registers & DSP / BRAM & Timing/evidence \\
\midrule
Tang fast probe & 696 LUT4; 216 DFF & 0 / 0 & 126.50 MHz; silicon \\
Tang PHSLK probe & 293 LUT4; 146 DFF & 0 / 0 & 200.40 MHz; silicon \\
Artix Lucas MAC leaf & 955 estimated LCs & 92 / -- & mapped in LUCAS spin \\
Whole Artix LUCAS spin & 5,073 estimated LCs & 120 / -- & 4.41 MHz; silicon \\
\bottomrule
\end{tabular}
\end{table*}
}

\newcommand{\LucasPinvProfileTable}{%
\begin{table}[ht]
\centering
\caption{Exhaustive software-oracle profile of the RTL-style scalar PINV
inverse over $\mathbb{Z}[\phi]/L_{521}$. Counts exclude acceptance.}
\label{tab:pinv-profile}
\footnotesize
\setlength{\tabcolsep}{3pt}
\begin{tabular}{lr}
\toprule
Measure & Oracle result \\
\midrule
Nonzero-norm elements & 270,400 \\
Zero-divisor elements & 1,041 \\
Unique nonzero norms & 520 \\
Busy cycles, min/max/mean & 3/23/16.66 \\
\bottomrule
\end{tabular}
\end{table}
}
% End auto-generated LUCAS publication data
